\section{Data}
\label{sec:data}

Downscaling requires fine-resolution predictors and coarse labels.
Features are drawn from several sources (Table~\ref{tab:data-sources}): Sentinel-2 reflectance, Xavier climate covariates, and MapBiomas Solo soil properties at each soybean pixel, a MapBiomas soybean mask, and IBGE PAM productivity labels at the municipality level.

\begin{table}[t]
\centering
\caption{Summary of data sources used in the experiments.}
\label{tab:data-sources}
\footnotesize
\setlength{\tabcolsep}{3.5pt}
\begin{tabular}{@{}llp{2.35cm}ccp{2.35cm}@{}}
\toprule
Category & Variables & Unit & Related properties & Spatial res. & Source \\
\midrule
Satellite imagery
  & \makecell[l]{B2, B3, B4, B5, B6,\\B7, B8, B8A, B11, B12}
  & ---
  & Vegetation reflectance
  & 30\,m
  & Sentinel-2 L2A \\
\midrule
\multirow{6}{*}{Climate}
  & Cumulative precipitation
  & mm
  & Water availability
  & \multirow{6}{*}{$0.1^\circ$}
  & \multirow{6}{*}{\makecell[c]{Xavier\\(BR-DWGD)}} \\
\cmidrule(lr){2-4}
  & Max.\ dry-day streak
  & days
  & Drought stress
  &
  & \\
\cmidrule(lr){2-4}
  & Cumulative ETo
  & mm
  & Atmospheric demand
  &
  & \\
\cmidrule(lr){2-4}
  & Cumulative Rs
  & MJ\,m$^{-2}$
  & Solar radiation
  &
  & \\
\cmidrule(lr){2-4}
  & Cumulative Tmax
  & $^\circ$C$\cdot$day
  & Heat accumulation
  &
  & \\
\cmidrule(lr){2-4}
  & Cumulative Tmin
  & $^\circ$C$\cdot$day
  & Cold accumulation
  &
  & \\
\midrule
\multirow{4}{*}{Soil}
  & Clay (0--30\,cm)
  & \%
  & Texture
  & \multirow{4}{*}{30\,m}
  & \multirow{4}{*}{\makecell[c]{MapBiomas Solo\\Col.\ 3}} \\
\cmidrule(lr){2-4}
  & Silt (0--30\,cm)
  & \%
  & Texture
  &
  & \\
\cmidrule(lr){2-4}
  & Sand (0--30\,cm)
  & \%
  & Texture
  &
  & \\
\cmidrule(lr){2-4}
  & SOC stock (0--30\,cm)
  & \si{t/ha}
  & Organic carbon
  &
  & \\
\midrule
Land cover
  & Soybean mask (class 39)
  & ---
  & Crop mask
  & 30\,m
  & MapBiomas Col.\ 10 \\
\midrule
Labels
  & Soybean productivity
  & \si{t/ha}
  & Municipal supervision
  & Municipality
  & IBGE PAM \\
\bottomrule
\end{tabular}
\end{table}

\subsection{Study area and labels}
The study covers soybean-producing municipalities in the state of Paran\'{a}, Brazil.
Supervision uses municipal soybean productivity in tonnes per hectare (\si{t/ha}) from IBGE PAM, with municipality codes, harvest year, and a predefined validation and test partition within the holdout year.
No field- or pixel-level yield labels are used.
We retain municipality--year pairs whose Sentinel-2 coverage ratio relative to the reported planted area lies in $[0.8, 1.2]$, so that satellite coverage is roughly consistent with official area statistics.

\subsection{Imagery and climate covariates}
Predictors are per-pixel Sentinel-2 spectral time series, organized by municipality and growing season.
Each pixel sequence comprises ten spectral bands (B2, B3, B4, B5, B6, B7, B8, B8A, B11, B12) and a day-of-year index used for positional encoding.
Climate covariates come from the Xavier daily weather grid~\citep{xavier2022} and are temporally aligned to each Sentinel-2 observation date: cumulative precipitation from 1~October through the observation date; the maximum consecutive dry-day streak between successive Sentinel-2 dates (dry day: precipitation $<1$\,mm); and season-to-date cumulative reference evapotranspiration (ETo), solar radiation (Rs), maximum temperature (Tmax), and minimum temperature (Tmin).
All climate channels are sampled at the Xavier cell containing each pixel centroid and scaled before concatenation with normalized reflectance.
The proposed model uses all six Xavier climate channels together with the ten spectral bands (sixteen features per timestep).
Sequences are truncated or padded to a fixed length of $T=45$ timesteps.
Soybean pixels are selected with the MapBiomas Collection~10 land-cover product (class~39), which defines the intramunicipal support of the downscaled maps.
An ablation omits all climate covariates and uses the ten spectral bands alone.

\subsection{Train / validation / test split}
Splits follow a year-holdout design on municipalities.
All municipality--year observations outside the holdout year are used for training; municipalities in the holdout year are partitioned into validation and test sets.
The holdout year is \textbf{2021}.
After restricting to harvest years $\{2020,2021,2022,2023,2024\}$, the training set contains 997 municipality--years, while validation and test contain 132 and 123 municipalities in 2021, respectively.
Training target statistics used for loss normalization are mean $3.17$\,\si{t/ha} and standard deviation $0.93$\,\si{t/ha}.
Quantitative evaluation uses municipal aggregates of the pixel predictions; intramunicipal structure is examined qualitatively through pixel maps.
